\documentclass[11pt]{article}
\usepackage[a4paper,margin=1in]{geometry}
\usepackage[T1]{fontenc}
\usepackage{lmodern}
\usepackage{microtype}
\usepackage{parskip}

\title{Traditional Event Representation Baseline Analysis}
\author{}
\date{}

\begin{document}
\sloppy
\maketitle

\paragraph{Saturated Recognition On N-MNIST} The N-MNIST classification results show that the task is nearly saturated under the unified ResNet18 classifier: four of the five traditional representations exceed 98.0\% test accuracy, and Voxel Grid reaches the best accuracy of 99.08\% with the lowest test loss of 0.034060. The spread between the best method, Voxel Grid at 99.08\%, and the weakest method, Binary Event Image at 95.06\%, is 4.02 percentage points, indicating that temporal binning still matters even on this comparatively simple recognition benchmark. At the same time, the ranking is not purely determined by training time: Timestamp Image is the fastest run at 786.48 seconds and still reaches 98.05\%, whereas Event Frame takes 2365.40 seconds for 98.28\%. This suggests that, for N-MNIST, compact timestamp-based encoding offers a strong efficiency--accuracy trade-off, while Voxel Grid provides the best absolute accuracy at moderate additional cost.

\paragraph{Temporal Encoding Matters More On N-Caltech101} N-Caltech101 exposes a clearer gap among traditional representations than N-MNIST. Timestamp Image obtains the best test accuracy, 52.9277\%, and is also the fastest method at 793.65 seconds, followed by Time Surface with 51.7796\% at 838.09 seconds. In contrast, Binary Event Image obtains the lowest accuracy, 48.6797\%, despite a longer runtime of 956.50 seconds. The gap between the best and weakest methods is 4.248 percentage points, close to the spread observed on N-MNIST, but the absolute accuracy is much lower, suggesting that N-Caltech101 is more sensitive to the representational loss caused by collapsing event timing. Voxel Grid reaches a validation accuracy tied with Timestamp Image at 50.0717\%, but its test accuracy drops to 49.6556\%, so the validation advantage does not translate into the best held-out performance in this run.

\paragraph{Traditional Baselines Remain Competitive On MVSEC} Under the unified MVSEC optical-flow protocol, all 11 methods use the same training split, evaluation split, window size of 200, stride of 200, shared EVFlowNet-like decoder, and early stopping with patience 10, minimum delta 0.001, and 1000 block-random outdoor validation windows. Within this controlled setting, Timestamp Image obtains the best overall AEE of 2.791623 and the lowest outlier rate of 36.608703\%, outperforming the best learning-based method, EST, which obtains an AEE of 2.865447 and an outlier rate of 37.039154\%. The AEE difference is 0.073824, corresponding to a 2.576\% relative reduction against EST. This result should not be interpreted as a universal dominance claim over learning-based methods, because the comparison uses a shared benchmark decoder rather than each paper's original full training stack. It does, however, show that a simple timestamp-preserving traditional representation can be highly competitive when evaluated under the same downstream optical-flow protocol.

\paragraph{Representation Complexity Does Not Guarantee Better Accuracy} Across the MVSEC comparison, channel dimensionality is not monotonically associated with better AEE. Timestamp Image uses only 2 channels and achieves the best AEE, while EST uses 18 channels and ranks second with AEE 2.865447. Voxel Grid uses 10 channels and ranks third with AEE 2.900117, whereas GET uses 24 channels but obtains AEE 2.961931. The 2-channel Event Frame and the 2-channel Event Pre-training representation are nearly indistinguishable in AEE, 2.965373 versus 2.965345, suggesting that, in this particular unified decoder setting, additional learned or structured channels do not automatically translate into stronger optical-flow estimates. The more conservative conclusion is that temporal information placement appears more important than raw channel count for this protocol.

\paragraph{Early Stopping Reveals Different Optimization Profiles} The early-stopping behavior also separates the representations. On MVSEC, Timestamp Image and EST both select the first epoch as the best validation checkpoint and stop after 11 epochs, whereas Voxel Grid reaches its best validation AEE at epoch 13 and stops after 23 epochs. Binary Event Image, Event Frame, Time Surface, Event Pre-training, and ERGO all converge to their selected best validation checkpoints around epoch 10 and stop after 20 epochs. Thus, the strongest MVSEC method in this run is also one of the fastest to validate, while Voxel Grid requires a longer optimization trajectory to reach a competitive but lower-ranked final AEE. This pattern supports using early-stopped validation curves in the report rather than only reporting final test metrics, because methods with similar final AEE can have noticeably different training dynamics.

\paragraph{Cross-Task Interpretation} The three evaluated settings point to a consistent but limited conclusion: traditional representations remain meaningful baselines and should not be treated as trivial lower bounds. Voxel Grid is strongest on N-MNIST with 99.08\% accuracy, while Timestamp Image is strongest on N-Caltech101 with 52.9277\% accuracy and on MVSEC with AEE 2.791623. The shift from Voxel Grid on N-MNIST to Timestamp Image on N-Caltech101 and MVSEC suggests that the best traditional representation is task-dependent rather than universally fixed. Because GEN1 detection was not included in the completed results, no detection-side conclusion should be drawn from the current data. The appropriate reporting claim is therefore that traditional baselines are competitive under the available classification and optical-flow protocols, with timestamp-preserving methods providing the most consistent cross-task performance among the completed experiments.

\end{document}
